% Generated by Listas/make_statements.py from pratica-lecture-08.md, item 25. Do not edit.
\textbf{Source.} Mas-Colell, Whinston \& Green (1995), \emph{Microeconomic Theory}, Ch. 13 Exercises, Exercise 13.C.4\^{}B, p.~475 {[}PDF p.~248{]}. Reproduced verbatim. The upper bound's accent is blurred in the scan; read as \(\bar{\theta}\), the notation of §13.C and of 13.C.1. Answer: none in the book; entries for 13.C.1, 13.C.2, 13.C.4 and 13.C.5 are in the MWG solutions manual (\texttt{Livros/}).



\begin{quote}

\textbf{13.C.4\^{}B} Reconsider the signaling model discussed in Section 13.C, now assuming that worker types are drawn from the interval \([\underline{\theta}, \bar{\theta}]\) with a density function \(f(\theta)\) that is strictly positive everywhere on this interval. Let the cost function be \(c(e, \theta) = (e^2/\theta)\). Derive the (unique) perfect Bayesian equilibrium.

\end{quote}



\textbf{Model of Section 13.C (p.~450 {[}PDF p.~236{]}), copied from the book because the exercise refers to it.}



\begin{quote}

Consider the following adaptation of the model discussed in Section 13.B. For simplicity, we restrict attention to the case of two types of workers with productivities \(\theta_H\) and \(\theta_L\), where \(\theta_H > \theta_L > 0\) and \(\lambda = \text{Prob}\,(\theta = \theta_H) \in (0, 1)\). The important extension of our previous model is that before entering the job market a worker can get some education, and the amount of education that a worker receives is observable. To make matters particularly stark, we assume that education does \emph{nothing} for a worker's productivity (see Exercise 13.C.2 for the case of productive signaling). The cost of obtaining education level \(e\) for a type \(\theta\) worker (the cost may be of either monetary or psychic origin) is given by the twice continuously differentiable function \(c(e, \theta)\), with \(c(0, \theta) = 0\), \(c_e(e, \theta) > 0\), \(c_{ee}(e, \theta) > 0\), \(c_\theta(e, \theta) < 0\) for all \(e > 0\), and \(c_{e\theta}(e, \theta) < 0\) (subscripts denote partial derivatives). Thus, both the cost and the marginal cost of education are assumed to be lower for high-ability workers; for example, the work required to obtain a degree might be easier for a high-ability individual. Letting \(u(w, e \mid \theta)\) denote the utility of a type \(\theta\) worker who chooses education level \(e\) and receives wage \(w\), we take \(u(w, e \mid \theta)\) to equal her wage less any educational costs incurred: \(u(w, e \mid \theta) = w - c(e, \theta)\). As in Section 13.B, a worker of type \(\theta\) can earn \(r(\theta)\) by working at home.

\end{quote}



\textbf{Optional extras.}

- Tadelis 15.4 ``Campaigning'' (p.~316): a second two-type signaling game with commitment by the rival, after item 10.

- Tadelis 16.1 (p.~335): uniqueness of beliefs in separating PBE of the item 10 game.

- Tadelis 16.2 ``Limit Pricing Revisited'' (p.~335): separation at monopoly quantities with a discount factor \(\delta\). As printed, the second quantity reads \(q_1^{1L} = 1.5\); by Claim 16.2 it is the high-cost type's, \(q_1^{1H} = 1.5\).

- Gibbons 4.5 (all pure PBE of Problem 4.3) and 4.6 (a signaling version of Fig. 4.1.1), p.~249.

- Muñoz-Garcia Ch. 9 Ex. 3--4 (pp.~394, 406: Spence with productivity- and utility-enhancing education), Ex. 6 (p.~432: cost uncertainty), Ex. 7 (p.~437: three types, three messages); Ch. 11 Ex. 1--3 (pp.~482, 487, 497) and Ex. 6 (p.~510). Pages from the PDF bookmarks with the registry offsets; not opened.



\begin{center}\rule{0.5\linewidth}{0.5pt}\end{center}

